\documentclass[11pt]{article}
\usepackage[margin=.78in,top=.65in,bottom=.65in]{geometry}
\usepackage{amsmath,amssymb,amsthm,hyperref}
\renewcommand{\normalsize}{\fontsize{10}{12}\selectfont}
\DeclareMathSizes{10}{10}{8}{6}
\hypersetup{colorlinks=true,linkcolor=blue,urlcolor=blue,citecolor=blue}
\setlength{\emergencystretch}{2em}
\newtheorem{theorem}{Theorem}
\newtheorem{lemma}[theorem]{Lemma}
\newcommand{\R}{\mathbb R}
\newcommand{\E}{\mathbb E}
\begin{document}
\begin{center}
{\Large\bfseries A variance bound for Davis's crowd-estimation polynomials}\\[.5em]
Research proof and Lean formalization\\
October 27, 2026
\end{center}
\begin{abstract}
We prove the rational-function bound proposed in Conjecture~3.1 of Davis's
paper on crowd estimation. The inequality holds strictly throughout its
stated domain. The proof rewrites the polynomial quotient as a variance,
derives an exact fourth-moment identity by finite detailed balance, and
uses a quadratic comparison for every integer degree at least~43.
The remaining~38 degrees are covered by exact positive-coefficient
certificates. The source's separate conjecture on critical points is
not addressed.
\end{abstract}

\section{Statement and the finite distribution}
For $i\ge0$ write $f_i(x)=\prod_{j=0}^{i-1}(x-j)$, with $f_0=1$. Define
\begin{align}
 P_d(x)&=\sum_{i=0}^d\frac{f_i(x)^2f_{d-i}(x)^2}{i!(d-i)!},\label{Pd}\\
 q_d(x)&=\frac{2x^2-2dx}{d+1}+\frac d2,
 \qquad R_d(x)=P_{d+1}(x)-q_d(x)P_d(x).\label{remainder}
\end{align}
Davis proves that $q_d$ is the polynomial quotient and $R_d$ the remainder
(Lemma~1.6 and equation~(2.4) in the arXiv version; Lemma~1.1 and
equation~(2.4) in the revised journal-style manuscript)~\cite{Davis}.
His Conjecture~3.1 is the following inequality with a weak conclusion.
\begin{theorem}\label{main}
For every integer $d\ge5$ and real $x\ge(d+\sqrt{d+2})/2$,
\[
 -\frac{R_d(x)}{P_d(x)}>
 \frac{199}{200}\,\frac{d-2}{2x+d-2}.
\]
\end{theorem}
The constant is the source's $0.995$. We make no claim that it is optimal.

Put $w_i=f_i(x)^2f_{d-i}(x)^2/(i!(d-i)!)$ and $Z=P_d(x)$.
These weights are nonnegative and $Z>0$ when $x>d/2$: the term
$i=\lfloor d/2\rfloor$ has only positive factors in both falling
factorials. Let expectations use probabilities $w_i/Z$, and set
\[
 D=d/2,\quad a=x-D,\quad u=i-D,\quad V=\E u^2,\quad M=\E u^4.
\]
Symmetry gives $\E i=D$. If $w'_j$ denotes the degree-$(d+1)$ weights,
then $(i+1)w'_{i+1}=(x-i)^2w_i$. Summing and using symmetry again yields
\begin{equation}\label{variance}
 \frac{P_{d+1}}{P_d}=\frac{2}{d+1}\E(x-i)^2
 =\frac{2}{d+1}(a^2+V),\qquad
 -\frac{R_d}{P_d}=\frac{d-4V}{2(d+1)}.
\end{equation}
Thus an upper bound for $V$ suffices.

\section{An exact fourth-moment identity}
For $0\le i\le d$ put
\[
 \lambda_i=(d-i)(x-i)^2,\qquad \mu_i=i(x-d+i)^2.
\]
Direct multiplication gives
$\lambda_iw_i=\mu_{i+1}w_{i+1}$ for $i<d$, and
$\lambda_d=\mu_0=0$. Hence, for any function $h$ on the integers,
\begin{equation}\label{balance}
 \sum_{i=0}^dw_i\{\lambda_i[h(i+1)-h(i)]
                 +\mu_i[h(i-1)-h(i)]\}=0.
\end{equation}
This is a finite index-shift identity. No division by $x-i$ is used,
so vanishing weights at integer inputs cause no exception.

\newpage
Take $h(i)=(i-D)^2$. Expansion gives
\[
 \lambda-\mu=-2u(2Da+a^2+u^2),\qquad
 \lambda+\mu=2(Da^2+Du^2+2au^2).
\]
Substituting into~\eqref{balance} and dividing by $2Z$ proves
\begin{equation}\label{moment}
 2M+BV=Da^2,\qquad B=4Da-D+2a^2-2a.
\end{equation}
Since $\E(u^2-V)^2\ge0$, we have $M\ge V^2$. Therefore
\begin{equation}\label{quadratic}
 2V^2+BV\le Da^2.
\end{equation}

\section{Degrees at least 43}
The source domain implies $a\ge\sqrt{d+2}/2>1$ and $4a^2\ge d+2$.
Define
\[
 T=\frac d4-\frac{199(d+1)(d-2)}{800(a+d-1)}
   =\frac{200da+d^2-d+398}{800(a+d-1)}>0.
\]
By~\eqref{variance}, the conclusion of Theorem~\ref{main} is precisely
$V<T$. Also $B>0$. The needed comparison is certified by
\begin{align*}
 2T^2+BT-Da^2&=\frac{N(d,a)}{320000(a+d-1)^2},\\
 N(d,a)&=800K(a+2d-2)\left(a^2-\frac{d+2}{4}\right)
              +aL(d)+K(201d^2+599d-402),\\
 K&=d^2-d+398,\\
 L(d)&=400(2d^4-105d^3+902d^2-1197d+1194).
\end{align*}
Every term in $N$ is nonnegative, and the last two are positive.
Indeed $K>0$, and for $t=d-43\ge0$,
\[
 L(d)=800t^4+95600t^3+3818000t^2+51998400t+42755200>0.
\]
Thus $2T^2+BT>Da^2$. On $[0,\infty)$ the polynomial $2v^2+Bv$
is strictly increasing. Since $V\ge0$, inequality~\eqref{quadratic}
forces $V<T$. This proves the entire infinite range.

\section{Degrees 5 through 42}
Let $S_d=d!P_d$, an integer polynomial. To avoid fractions define
\begin{align*}
 H_d(x)={}&\{200(2x+d-2)[4x^2-4dx+d(d+1)]
                   -398(d+1)(d-2)\}S_d(x)\\
          &-400(2x+d-2)S_{d+1}(x).
\end{align*}
An exact identity is
\begin{equation}
 H_d=400(d+1)d!\left\{
    [(2x+d-2)q_d-(199/200)(d-2)]P_d-(2x+d-2)P_{d+1}\right\}.
\end{equation}
Since $P_d>0$ and $2x+d-2>0$, positivity of $H_d$ is equivalent to
the strict desired bound.

Put
\[
 m_d=1000d+\left\lfloor\sqrt{10^6(d+2)}\right\rfloor,
 \qquad r_d=m_d/2000\le(d+\sqrt{d+2})/2.
\]
Represent $H_d$ by its ascending integer coefficient list $A_d$;
trailing zeros are permitted. If its list length is $\ell_d$, set
\[
 B_d(y)=2000^{\ell_d-1}H_d(y+r_d).
\]
For each of the~38 integers $d=5,\ldots,42$, the constant coefficient
of $B_d$ is positive and every coefficient is nonnegative.
These are exact integer certificates, checked from the defining falling
products. Consequently $B_d(y)>0$ for $y\ge0$, hence $H_d(x)>0$
for every $x\ge(d+\sqrt{d+2})/2$. Together with the preceding section
this proves Theorem~\ref{main}.

\newpage
\section{Formal certificate and scope}
The Lean development uses the literal finite sum~\eqref{Pd} and
difference~\eqref{remainder}. It proves positivity of $P_d$, reflection
and the mean identity, detailed balance, the fourth-moment equation,
$M\ge V^2$, and the uniform scalar comparison. There are no residual
moment or polynomial-identity hypotheses in the source conclusion.
The endpoint names are
\begin{center}
\texttt{Davis.source\_strict\_bound},\qquad
\texttt{Davis.source\_conjecture\_319187}.
\end{center}
The finite checker constructs all coefficients internally from the
source formulas. Polynomial addition, multiplication, scalar multiplication,
and the affine translation are implemented on integer lists, with their
evaluation identities proved in Lean. Binomial coefficients use the
factorial implementation whose equality with the usual binomial coefficient
is proved in mathlib. The~38 positivity checks use kernel reduction
of integer arithmetic. They do not invoke an untrusted external numerical
oracle. The rational shift is bounded using the proved inequality
$\lfloor\sqrt n\rfloor^2\le n$.

The certificate has 100 declarations in 19 modules. A fresh isolated replay
with Lean~4.34.1 and mathlib revision \texttt{d13f23b7} passes with
warnings as errors. Every declaration type and axiom report was inspected;
only \texttt{propext}, \texttt{Classical.choice} and \texttt{Quot.sound}
occur.

As supplementary checks, a pure-integer Python implementation and a
separate symbolic rational-polynomial implementation produce positive
certificates for all~38 finite degrees. Exact rational evaluation verifies
the weight, quotient and moment identities at~174 inputs, including~45
cases with zero weights. These checks are independent of the universal
mathematical argument and are not used to replace the Lean proof.

The source's established half-unit bound, polynomial quotient,
conditional crowd-estimation consequences, and numerical observations
are credited to Davis. Detailed balance and the variance inequality are
standard tools; the proposed contribution is their combination with the
explicit uniform comparison and finite certificates to resolve the exact
source conjecture. We do not prove the distinct critical-point
Conjecture~3.9 of the arXiv version (Conjecture~3.2 in the revised
manuscript), identify the optimal constant, or claim new general
moment theory.

The prior-art investigation is bounded. The current arXiv entry has
only version1, the author's revised 2010 manuscript retains Conjecture~3.1,
and targeted searches located no resolution. The exact OpenAlex record
has no indexed citers; this count is not evidence that no unindexed or
unpublished resolution exists. Subscription indexes and a final paginated
publisher copy were not available. No absolute priority claim is made.

\begin{thebibliography}{9}
\bibitem{Davis} D. M. Davis, \emph{Maximizing a combinatorial expression
arising from crowd estimation}, International Journal of Modern
Mathematics \textbf{5} (2010),129--148.
\href{https://arxiv.org/abs/0910.5725v1}{arXiv:0910.5725v1}.
Author-hosted revised journal-style manuscript (May 14, 2010):
\url{https://www.lehigh.edu/~dmd1/fredijmm.pdf}.
\end{thebibliography}
\end{document}
